\documentclass{article}
\usepackage[utf8]{inputenc}
\usepackage[english]{babel}
\usepackage{graphicx}
\usepackage[authoryear]{natbib}
\usepackage{titlesec}
\graphicspath{ {./images/} }
\newcommand{\sectionbreak}{\clearpage}

\newenvironment{changemargin}[2]{%
\begin{list}{}{%
\setlength{\topsep}{0pt}%
\setlength{\leftmargin}{#1}%
\setlength{\rightmargin}{#2}%
\setlength{\listparindent}{\parindent}%
\setlength{\itemindent}{\parindent}%
\setlength{\parsep}{\parskip}%
}%
\item[]}{\end{list}}

\title{BTP Report}
\author{Naman Jain, Pranshu Maheshwari, Saksham Tanwar}
\date{\today}

\begin{document}

\begin{titlepage}
    \begin{changemargin}{-2cm}{-2cm}
        \begin{center}
            \textbf{\Huge{B. TECH. PROJECT REPORT}\\}
            \textbf{\Huge{on}\\*[1cm]}
            \textbf{\Huge{Smartphone based Land Records Retrieval System}\\}
            \vspace{3cm}
            By\\
            \textbf{Naman Jain}\\
            \textbf{Pranshu Maheshwari}\\
            \textbf{Saksham Tanwar}\\

            \begin{figure}[htbp]
                \includegraphics[scale=0.5]{iiti_logo}
                \centering
            \end{figure}
            \vspace{3cm}
            \textbf{\large{DISCIPLINE OF COMPUTER SCIENCE AND ENGINEERING}}\\*[0.2cm]
            \textbf{\LARGE{INDIAN INSTITUTE OF TECHNOLOGY INDORE}}\\
            \today
        \end{center}
    \end{changemargin}
\end{titlepage}
\clearpage
\tableofcontents
\clearpage

% \section*{Sections in Report}
% \begin{enumerate}
%     \item Introduction
%     \begin{itemize}
%         \item Description of a generic application
%         \item Disadvantages of a centralized database layer
%         \item Advantages of using a blockchain solution
%         \item Introduction of our product, highlighting the importance of blockchain, describing overall architecture and various products used.
%     \end{itemize}
%     \item Introduction of essential Concepts of Blockchain
%     \begin{itemize}
%         \item Describe various offerings of blockchain in regards to our product
%         \item Existing applications related to our product
%     \end{itemize}
%     \item Architecture
%     \begin{itemize}
%         \item Details regarding the solution architecture
%         \item Identifying the various actors/users of the product and how they interact with the product
%         \item Various sub modules of the product and how they interact with one another
%     \end{itemize}
%     \item Deployment Details
%     \begin{itemize}
%         \item Details about how the product can be deployed
%     \end{itemize}
%     \item Conclusion
%     \begin{itemize}
%         \item Identify the problems which the product solved
%         \item Identify the problems still left in the solution proposed
%     \end{itemize}
%     \item References

% \end{enumerate}

\section{Introduction}
    \subsection{Generic Application}
        \paragraph{} A typical application can be broken down into three logical layers namely:-
        \begin{itemize}
            \item \textbf{Presentation Layer: }It is the layer responsible for display information and collect information from the user.
            \item \textbf{Business Logic Layer: }In this layer, information collected from presentation layer is used to perform the various calculations and operations needed to be performed by the application.
            \item \textbf{Data Layer: }This layer is responsible for storing and retrieving data to be used by other layers.
        \end{itemize}

        \paragraph{}Data Layer is typically implemented using a database, which is an organized collection of data. A database can be of two types : Centralized and Distributed.
        
    \subsection{Disadvantages of a centralized database}
        \paragraph{}
        Most applications require data in some form or another and need to store the said data for future use, this is where databases come into picture.
        A database is an organized collection of data, generally stored and accessed electronically from a computer system.
        Access to this data is usually provided by a "database management system" (DBMS) consisting of an integrated set of computer software that allows users to interact with one databases and provides access to all of the data contained in the database.
        DBMS's provide various functions that allow management of a database and its data such as Data definition which defines the database model i.e. the logical structure of the data model, Updating of actual data which includes insertion, modification, and deletion of the actual data, Retrieval of actual data which may or may not include data processing and Administration which includes allowing access to various users, enforcing data security and various other checks.
        \cite{wiki}
        A database model is a type of data model that determines the logical structure of a database and fundamentally determines in which manner data can be stored, organized and manipulated. The most popular example of a database model is the relational model, which uses a table-based format. The ACID (Atomicity, Consistency, Isolation and Durability) database design is one of oldest and most important database design. A relational database that fails to follow the ACID model cannot be considered reliable. 
        \paragraph{}
        Databases can categorized in many categories depending on how they store data, where it is stored and the additional functionality the DBMS can provide on the stored data. Broadly databases can categorized into two major categories centralized and distributed. They both have their advantages and disadvantages.
        A centralized database is a database that is located, stored, and maintained in a single location and has a single DBMS but can be accessed by the users distributed in the network.
        A distributed database system consists of a collection of local databases,  geographically located in different points (nodes of a network of computers) and  logically related by functional relations so that they can be viewed globally as a single database \cite{distributeddatabase}.
        There are two types of distributed databases homogeneous and heterogeneous. In a homogeneous distributed database all the locations store the database identically i.e. all locations have the same management system and schema. Whereas in heterogeneous distributed database different locations can have different management software, schemas.
        For a database management system to be distributed, it should be fully compliant with the twelve rules introduced by C.J. Date in 1987 \cite{distributeddbms}: local autonomy; the absence of a dependency from a central location; continuous operation; location independent; fragmentation independent; replication  independent; distributed query processing; distributed transaction management; hardware independent; operating system independent; independent of communication  infrastructure; independent of database management system.
        
        \paragraph{}
        If a centralized database is used to store the data then the whole system is prone to Single Point of Failure, if for some reason the database fails the whole system will fail. As all the data is stored in a single location, if there are no  precautionary measures taken then a hardware failure can lead to complete data loss. The databases are also vulnerable to cyberattacks which can cause data loss, data leaks, data inconsistency and many other vulnerability. This makes centralized database dubious with very low reliability. But at the same time, centralized database ensures data consistency and easy management in compliance with ACID design \cite{centralizeddistributeddatabases}.
        
        \paragraph{}
        A major advantage of distributed database is that by sharing a database across  multiple nodes can obtain a storage space extension and also can benefit from multiple processing resources. A distributed database system is robust to failure to some extent. Hence, it is reliable when compared to a  centralized database system. It is also more robust compared to a centralized database as it doesn't have a single point of failure, if a node fails another node or group of nodes can provide the necessary data. But to make this possible complex software's are required which incur additional costs and processing overheads. Marinating data integrity is difficult and hence minimum redundancy and ACID properties are more relaxed than compared to a centralized database \cite{centralizeddistributeddatabases}. Distributed environment also faces problems such as fragmentation and data replication. A data fragment constitutes some subset of the original database. A data replica constitutes some copy of the whole or part of the original database. The fragmentation and the replication can be combined: a relationship can be partitioned into several pieces and can have multiple replicas of each fragment \cite{distributedsystems}.
        
    \subsection{Advantages of using Blockchain}
        % general introduction about blockchain benefits.
        Blockchain is a type of decentralized database where the data is stored in blocks and each block is linked to the previous block using its cryptographic hash, thus forming a chain. As the blockchain grows in size, it becomes more and more difficult to tamper the data in these blocks. Blockchain based databases are more trustworthy, reliable and secure than traditional databases.

        \paragraph{}
        % decentralization benefits
        A blockchain based system is decentralized. Hence, a centralized authority is not necessary. Decentralized systems are also resilient to single point of failure. A blockchain network comprises of various nodes, each maintaining its own copy of database. Hence, all the data is replicated, shared and synchronized across all the nodes in the network. This makes the data almost tamper proof as it will be very costly and thus highly unlikely for a malicious person to change the data across the whole network. No node can directly write data to the blockchain as it can make the data inconsistent, to avoid this the blockchain network uses a consensus algorithm which helps all the nodes in the network to come to a consensus and commit the data in the blockchain. This helps to keep all the nodes in the network to get synchronized with each other. 

        \paragraph{}
        % Immutability and tracability benefits.
        The data that is committed to the blockchain is immutable. After some data is inserted into the blockchain, it is very difficult to delete or alter it. This is realized by including, in each block, the cryptographic hash of its previous block. Because of this immutability, a blockchain database only supports create and read operations in contrast to traditional databases which supports create, read, update and delete (CRUD) operations. To update the database state, a subsequent write is required. Old writes persist in the blockchain forever, making it convenient to trace how the blockchain database state is updated over time. This makes it possible to trace the actions that resulted in a particular database state.

        \paragraph{}
        % transparency
        Another benefit of blockchain is that it provides transparency to the data stored in it. This is achieved by replicating the blockchain data among multiple nodes present in the blockchain network. Nodes in the blockchain network not only can view the data already committed to the blockchain but can also participate in the process of validating the data to be committed in the blockchain. The transparency also helps users to cross verify some data against the blockchain.  

        \paragraph{}
        % single public ledger
        A blockchain ledger is logically a single public ledger. Since every node updates its copy of blockchain ledger in sync with each other and after consensus is reached, the whole network logically act as one single public ledger. This removes the complications present in traditional systems where different participants or organizations maintains multiple ledgers which needs to be reconcile and synchronized time to time.

        \paragraph{}
        % Reduced Cost
        Systems that are deployed on public blockchain network are more cost efficient as compared to the traditional systems. This is because the system utilizes the processing power and resources of a large number of nodes that are already connected to the blockchain network, thus significantly reducing the cost needed for setting up and maintaining centralized servers present in traditional systems.

        \paragraph{}
        % Security
        Blockchain based systems are more secure and resilient to cyber attacks, which can cause data loss, data leaks, data inconsistency and many other vulnerability, than traditional sytems. Whenever a new block is introduced in the blockchain, its hash value is calculated which also includes the hash of its previous block. If a malicious user fraudulently tries to tamper the data inside a block, not only its hash value changes, but the hashes of the following blocks get changed too. This, along with the fact, that the data is stored in multiple nodes present in the blockchain network, makes it very difficult to tamper the data present in the blockchain.
        
        \paragraph{}
        % Benefits of blockchain in our solution

    \subsection{IPFS}

        \paragraph{}
        The Interplanetary File System (IPFS) \cite{IPFS} is a peer to peer distributed file system for sharing and storing data. IPFS provides a high throughput storage system where each file is uniquely identified by a cryptographic hash. In IPFS, multiple users store a piece of each others' data and actively participate in making it available in the network. In contrast to HTTP which uses location based addressing, IPFS uses content based addressing. In content based addressing, user request data using its content identifier instead of the location where it is stored.

        \paragraph{}
        IPFS uses Distributed Hash Table (DHT) \cite{DHT} to enable routing and discovery of peers and content on the network. The hash table is split across all the participating nodes in the IPFS. These nodes coordinate with each other to enable efficient lookup and retrieval of data on the network. IPFS stores data using a data structure called Merkle DAGS \cite{IPFS}. Each node in a Merkle DAG stores the cryptographic hash of each of its children. Thus, altering a node's data changes the hash value of the node as well as all its ancestors in the DAG. This makes the data stored on the IPFS immutable.
        
    \subsection{Motivation}
        \paragraph{}
        In India, currently if person is asked to prove if he/she is the owner of a piece of land then all they can show for it is a sale deed which just proves that the person was the owner at a particular time, but that person cannot prove that he/she is still the owner of that land. For proving the same he/she has to go to various government offices and collect various document showing that no sale deed has been registered for that land after the one which the person has. When a person wishes to buy land in India, they have to be very careful and perform various checks such as
        \begin{itemize}
            \item Check if the deed title is in the name of the seller and if he/she has the full right to sell it
            \item Procure a Encumbrance Certificate from sub-registrar's office where the deed is registered which declares that the land is free of any legal hassle and unpaid dues
            \item Check if the Property tax and other bills are paid in full and the seller has the respective original receipts
            \item Check if the loan on the land has been completely repaid
        \end{itemize}
        If a buyer is lethargic in verifying these documents he/she can be easily duped into buying a disputed land parcel or it may even happen that the land in question did not even legally belong to the said seller, that is the seller was not the genuine owner of that land provided fake documents. As these documents are mostly maintained offline as hard copies these cases occur often.

        \paragraph{}
        Getting all the information such as the sale deeds and ownership history for a land parcel is very cumbersome, a simple solution to this problem is a smartphone application with which users/potential buyers can get land ownership history and related documents for a particular land parcel. These documents and data can be stored in a database by the government and will be updated every time a new sale deed is registered i.e. owner for a land parcel changes. This makes the process of acquiring the documents in question very easy and hassle free.

        \paragraph{}
        Indian government launched Digital India Land Record Modernization Programme in 2008. Main aim of this programme is to create a system of updated land records, automated and automatic mutation, integration between textual and spatial records, inter-connectivity between revenue and registration, to replace the present deeds registration and presumptive title system with that of conclusive titling with title guarantee \cite{dilrmp}. In addition to this, Indian government also plans on storing the digital records on a private blockchain and provide a system where citizens can view the ownership details and complete history of the property before going in for the purchase \cite{blockchaingovin}.  

    \subsection{Proposed Solution}
        To takle this problem of accessing land records, a smartphone based solution is proposed. Users can request land records for a particular land parcel by visiting the land parcel or marking it on the map. Land records are stored on a private blockchain which provides methods to access the data. For help in migrating old records to the new system, scanned documents can be uploaded. These documents are stored in a private IPFS cluster with their Content Identifier (CID) is stored on the blockchain. A verification portal is also part of the solution where users while accessing records will be given verification code which users can use to generate a digitally signed PDF for verifying the land records. 
        
    \subsection{Organisation of Dissertation}
        A brief introduction to our proposed solution and the motivation behind it was discussed in chapter one. Chapter one also provided the generic application architecture and contrasted the centralized databases with blockchain databases. Chapter two provides an introduction to the various offerings of the blockchain technology that is makes it useful for the proposed solution. Chapter two also provides a summary of related works. Chapter three describes the proposed solution's architecture including the various actors/users of the solution, various submodules as part of the solution and how they interact with one another. Chapter four provides details regarding the deployment of the proposed solution. Chapter five concludes with the main contributions of this work and proposes direction for future work.
        
\section{Advantages of Decentralized Solutions}

    \subsection{Blockchain}
        \paragraph{}
        \cite{THAKUR2020101940} discusses various complications present in the current land record management system in India. The existing land records in India are not administered properly due to which they often are not consistent with the ground truth. Various departments, at the village or District level, are involved in maintaining ownership history of land. Often, the data present with different departments are not in sync with each other which results in discrepancies. Such discrepancies often result in conflicts in land ownership, double spending problem, etc. The data is also susceptible to alteration by corrupt officials.

        \paragraph{}
        Blockchain provides a potential solution to various complications present in the current land record management system. Being a distributed database, Blockchain based land record system is resilient to single point of failures. Data is replicated across multiple nodes in the Blockchain network. Hence, it prevents any undesirable data loss as was seen during Haiti 2010 Earthquake \cite{haiti}.
        
        \paragraph{}
        A Blockchain ledger contains transactions that are endorsed by appropriate authorities. These transactions, being committed in a chronological order, provides a verifiable audit trail for legal investigations \cite{Sahlin2018BlockchainIA}. Once committed to the Blockchain, the transactions cannot be modified or deleted, hence providing a trusted source of information.

        \paragraph{}
        Blockchain is a shared ledger. Each participant has a copy of the ledger. With necessary permissions, a participant can commit data into the ledger which, after getting validated by other nodes, gets reflected across the entire Blockchain network. This removes the complications present in the current land record management systems where different departments maintain land data which needs to be reconcile and synchronized regularly.
    
    \subsection{IPFS}
        \paragraph{}
        IPFS, being distributed, is resilient to single point of failure. It is more secure to DDos attacks than centralized system. Multiple nodes can simultaneously participate in distributing a particular file resulting in better throughput. The data retrieved from IPFS can be verified by calculating its hash. Hence, data stored on IPFS is self-verifiable. Data can be pinned on multiple nodes to ensure it persists on IPFS. This makes IPFS resilient to data loss.

    \subsection{Related Works}
        \paragraph{}
        To tackle the problem of accessing and verifying the land ownership records, this paper proposes a smartphone based solution. Using which users can request land ownership records for a particular land parcel by taking the advantage of the GPS module available in smartphones. These land ownership records are stored in a permissioned blockchain and to make it backwards compatible i.e. migrating old records to the new system, scanned copies of the documents can be stored in permissioned IPFS cluster with their Content Identifier (CID) stored on the same blockchain. Now whenever users want the ownership history of a land parcel they use the smartphone application, mark the piece of land and after successful completion of the process they get a PDF report with the ownership history of that said land parcel. This PDF report also contains verification codes for each ownership record they hold, these verification codes can then be used on a verification portal to generate a digitally signed PDF verifying that the given information is true. This solution hence provides a trustful and hassle free mechanism for people to access and verify land ownership records.

    \subsection{Organisation}
        \paragraph{}
        \cite{ALAM2020} proposes a three phase project for adoption of Blockchain based solution for land titling problem in Bangladesh and provides an implementation of a prototype model using Ethereum based blockchain.
        
    \clearpage
    \bibliographystyle{unsrt}
    \bibliography{citation}
\end{document}
